%% Datasheet for Datasets (Gebru et al., 2021)
%% ADS Cross-National Educational Benchmark (unified_v4)
%% Include via %% Datasheet for Datasets (Gebru et al., 2021)
%% ADS Cross-National Educational Benchmark (unified_v4)
%% Include via %% Datasheet for Datasets (Gebru et al., 2021)
%% ADS Cross-National Educational Benchmark (unified_v4)
%% Include via %% Datasheet for Datasets (Gebru et al., 2021)
%% ADS Cross-National Educational Benchmark (unified_v4)
%% Include via \input{datasheet} from main paper appendix.

\section{Datasheet for Datasets}
\label{sec:datasheet}

This datasheet follows the template of Gebru et al.~\cite{gebru2021datasheets}
and documents the ADS cross-national educational benchmark released with this paper.

% =====================================================================
\subsection{Motivation}
% =====================================================================

\paragraph{Why was the dataset created?}
No existing benchmark combines multi-institutional course artifacts with structured labor-market taxonomies and standardized multi-stakeholder evaluation tasks.
Learner-interaction datasets (e.g., KDD Cup 2010, EdNet, OULAD) focus on knowledge tracing at a single institution; MOOC-catalog datasets (MOOCCube, MOOCCubeX) cover multiple institutions but lack labor-market integration.
ADS was created to fill this gap: a cross-national, multi-objective educational benchmark for evaluating trade-offs among learner, institutional, and labor-market stakeholders.

\paragraph{Who created the dataset and on behalf of which entity?}
The dataset was created by the paper authors as an academic research contribution.

\paragraph{Who funded the creation of the dataset?}
No external funding was received for the creation of this dataset.
All data sources are publicly available and were collected using public APIs and catalogs.

% =====================================================================
\subsection{Composition}
% =====================================================================

\paragraph{What do the instances represent?}
The benchmark contains three types of primary text artifacts and two types of structured records:
\begin{itemize}
  \item \textbf{Courses} (\NumCourses): educational course descriptions from \NumUniversities{} universities across \NumCountries{} countries (USA, Sweden, UK).
  \item \textbf{O*NET occupation profiles} (\NumOccupations): occupation descriptions including required skills, tasks, and technologies from the US Department of Labor.
  \item \textbf{Mission statements} (\NumMissionsExpanded): university mission and vision texts.
  \item \textbf{College Scorecard programs} (\NumScorecardPrograms): program-level outcome descriptors with median earnings data (US DOE).
  \item \textbf{Structured attribute/link records} (\NumAttributeRecords): O*NET skills, tasks, technologies, and CIP-SOC crosswalk links.
\end{itemize}

\paragraph{How many instances are there?}
The benchmark contains \NumPrimaryArtifacts{} primary text artifacts and \NumProcessedRows{} total processed rows.
Course counts by university:
MIT OCW (\NumMIT),
UC Berkeley (\NumUCB),
Stanford (\NumStanford),
UIUC (\NumUIUC),
Cornell (\NumCornell),
KTH Sweden (\NumKTH),
Edinburgh UK (\NumEdinburgh).
US universities contribute \NumUSCourses{} courses; international universities contribute \NumIntlCourses{} courses.

\paragraph{What data does each instance consist of?}
Each \emph{course} record contains: \texttt{institution}, \texttt{course\_id}, \texttt{title}, \texttt{description}, \texttt{department}, and \texttt{level} (where available from the source catalog).
Each \emph{occupation} record contains: \texttt{occupation\_id}, \texttt{title}, \texttt{description}, lists of required skills, tasks, and technologies.
Each \emph{mission} record contains: \texttt{institution}, \texttt{statement\_text}, and \texttt{source\_url}.

\paragraph{Is there a label or target associated with each instance?}
No. The benchmark is designed for multi-objective evaluation, not supervised prediction. Evaluation is performed against the Pareto front and benchmark metrics (hypervolume, spacing, coverage).

\paragraph{Is any information missing from individual instances?}
Some fields are unavailable for certain sources: course \texttt{level} is not available from all university catalogs; \texttt{department} naming conventions vary across institutions.
International course descriptions (KTH, Edinburgh) are in English but may reflect translation artifacts from the original Swedish or course-catalog conventions.

\paragraph{Are there any errors, sources of noise, or redundancies?}
Near-duplicate courses (e.g., cross-listed courses) were removed during deduplication. Some course descriptions are short or formulaic (e.g., independent study entries); a minimum text-length filter was applied. O*NET profiles are professionally curated by the US DOL and contain minimal noise.

\paragraph{Does the dataset contain data that might be considered confidential or sensitive?}
No. All data is drawn from public catalogs, public APIs, and public-domain government sources. A basic PII screening step was applied to exclude any inadvertently included names, emails, or phone numbers.

% =====================================================================
\subsection{Collection Process}
% =====================================================================

\paragraph{How was the data associated with each instance acquired?}
Data was acquired from publicly accessible sources:
\begin{itemize}
  \item \textbf{MIT}: MIT OpenCourseWare (OCW) public catalog.
  \item \textbf{UC Berkeley}: UC Berkeley online course catalog.
  \item \textbf{Stanford}: Stanford ExploreCourses public API.
  \item \textbf{UIUC}: UIUC Course Information Suite (CIS) public API.
  \item \textbf{Cornell}: Cornell Class Roster public API.
  \item \textbf{KTH}: KTH KOPPS course syllabus public API (Sweden).
  \item \textbf{Edinburgh}: University of Edinburgh Degree Regulations and Programmes of Study (DRPS) public catalog (UK).
  \item \textbf{O*NET}: US Department of Labor O*NET public database and API.
  \item \textbf{College Scorecard}: US Department of Education College Scorecard public API.
\end{itemize}

\paragraph{What was the collection period?}
Data was collected between 2024 and 2025. O*NET time-travel versions (2022--2025) are included for market-drift analysis.

\paragraph{Were any ethical review processes conducted?}
All sources are public and open-access. No human subjects are involved; no IRB review was required. No personally identifiable information about students, instructors, or workers is collected or included.

\paragraph{Was consent obtained?}
Not applicable. All data consists of publicly published institutional artifacts (course catalogs, government databases).

% =====================================================================
\subsection{Preprocessing, Cleaning, and Labeling}
% =====================================================================

\paragraph{What preprocessing and cleaning was done?}
\begin{enumerate}
  \item \textbf{Schema validation}: all records validated against a required schema (\texttt{institution}, \texttt{course\_id}, \texttt{title} are mandatory fields).
  \item \textbf{Deduplication}: near-identical artifacts within each source were removed.
  \item \textbf{Minimum text length}: artifacts with extremely short or empty descriptions were filtered out.
  \item \textbf{PII screening}: basic screening to exclude names, email addresses, and phone numbers.
  \item \textbf{Provenance tracking}: for every artifact, we record a stable identifier, source URL, snapshot timestamp, content hash of the raw text, and parsing/extraction rules.
\end{enumerate}

\paragraph{Was the raw data saved in addition to the preprocessed data?}
Yes. Provenance manifests link each processed artifact to its raw source. Content hashes allow verification of data integrity.

\paragraph{Is the software used to preprocess the data available?}
Yes. Ingestion scripts and validation code are included in the released repository.

% =====================================================================
\subsection{Uses}
% =====================================================================

\paragraph{What tasks is the dataset intended for?}
The benchmark defines three evaluation tasks:
\begin{itemize}
  \item \textbf{Task A: Multi-Stakeholder Course Ranking.} Given course text and stakeholder target corpora, produce an ordered list of course IDs. Metrics: coverage, diversity, objective balance.
  \item \textbf{Task B: Stakeholder Trade-off Discovery.} Identify a set of courses representing distinct stakeholder trade-offs. Metrics: hypervolume, spacing, spread, Pareto precision, dominance resistance.
  \item \textbf{Task C: Robust Recommendation.} Produce a primary submission plus perturbation-run submissions. Metrics: bootstrap Jaccard, HV retention, rank stability, core stable fraction.
\end{itemize}

\paragraph{What tasks should the dataset \emph{not} be used for?}
The dataset is \emph{not} intended for:
\begin{itemize}
  \item Knowledge tracing or student performance prediction (no learner interaction data is included).
  \item Grade prediction or assessment modeling.
  \item Learner interaction modeling or tutoring system evaluation.
  \item Making enrollment decisions without human oversight.
\end{itemize}

\paragraph{Are there tasks for which the dataset should not be used?}
The dataset should not be used to rank individual students or to automate high-stakes enrollment/admissions decisions. It is a benchmark for evaluating multi-stakeholder trade-off methods, not an operational decision system.

% =====================================================================
\subsection{Distribution}
% =====================================================================

\paragraph{How will the dataset be distributed?}
Upon acceptance, the dataset will be distributed via:
\begin{itemize}
  \item \textbf{Zenodo} with a persistent DOI for long-term archival.
  \item \textbf{HuggingFace Datasets} for programmatic access.
\end{itemize}

\paragraph{When will the dataset be released?}
The dataset will be released upon paper acceptance.

\paragraph{What license will the dataset be distributed under?}
\textbf{CC BY-NC-SA 4.0}. This reflects the most restrictive component license (MIT OCW content is licensed under CC BY-NC-SA 4.0). O*NET data is public domain (US DOL). College Scorecard data is ODC-By (US DOE). University catalog data is from public sources.

\paragraph{Are there any fees or access restrictions?}
No. The dataset, evaluation script (\texttt{benchmark\_eval.py}, requiring only NumPy and the Python standard library), and ingestion code are freely available.

% =====================================================================
\subsection{Maintenance}
% =====================================================================

\paragraph{Who is supporting/hosting/maintaining the dataset?}
The dataset is maintained by the paper authors. Versioned snapshots are released; the current version is \texttt{unified\_v4}.

\paragraph{How can the owner/curator/manager of the dataset be contacted?}
Contact information will be provided upon de-anonymization.

\paragraph{Will the dataset be updated?}
The dataset follows a versioned-snapshot model. The current release (\texttt{unified\_v4}) includes courses from \NumUniversities{} universities. O*NET time-travel versions (2022--2025) are included to support market-drift analysis. Future snapshots may add universities, update course catalogs, or incorporate newer O*NET releases; each snapshot will receive a new version identifier and DOI.

\paragraph{Will older versions continue to be available?}
Yes. Each versioned snapshot is archived on Zenodo with its own DOI. Previous versions will remain accessible for reproducibility.

\paragraph{How will the dataset be retired?}
If a version is found to contain errors, a corrected version will be released and the errata will be documented. No version will be silently removed.
 from main paper appendix.

\section{Datasheet for Datasets}
\label{sec:datasheet}

This datasheet follows the template of Gebru et al.~\cite{gebru2021datasheets}
and documents the ADS cross-national educational benchmark released with this paper.

% =====================================================================
\subsection{Motivation}
% =====================================================================

\paragraph{Why was the dataset created?}
No existing benchmark combines multi-institutional course artifacts with structured labor-market taxonomies and standardized multi-stakeholder evaluation tasks.
Learner-interaction datasets (e.g., KDD Cup 2010, EdNet, OULAD) focus on knowledge tracing at a single institution; MOOC-catalog datasets (MOOCCube, MOOCCubeX) cover multiple institutions but lack labor-market integration.
ADS was created to fill this gap: a cross-national, multi-objective educational benchmark for evaluating trade-offs among learner, institutional, and labor-market stakeholders.

\paragraph{Who created the dataset and on behalf of which entity?}
The dataset was created by the paper authors as an academic research contribution.

\paragraph{Who funded the creation of the dataset?}
No external funding was received for the creation of this dataset.
All data sources are publicly available and were collected using public APIs and catalogs.

% =====================================================================
\subsection{Composition}
% =====================================================================

\paragraph{What do the instances represent?}
The benchmark contains three types of primary text artifacts and two types of structured records:
\begin{itemize}
  \item \textbf{Courses} (\NumCourses): educational course descriptions from \NumUniversities{} universities across \NumCountries{} countries (USA, Sweden, UK).
  \item \textbf{O*NET occupation profiles} (\NumOccupations): occupation descriptions including required skills, tasks, and technologies from the US Department of Labor.
  \item \textbf{Mission statements} (\NumMissionsExpanded): university mission and vision texts.
  \item \textbf{College Scorecard programs} (\NumScorecardPrograms): program-level outcome descriptors with median earnings data (US DOE).
  \item \textbf{Structured attribute/link records} (\NumAttributeRecords): O*NET skills, tasks, technologies, and CIP-SOC crosswalk links.
\end{itemize}

\paragraph{How many instances are there?}
The benchmark contains \NumPrimaryArtifacts{} primary text artifacts and \NumProcessedRows{} total processed rows.
Course counts by university:
MIT OCW (\NumMIT),
UC Berkeley (\NumUCB),
Stanford (\NumStanford),
UIUC (\NumUIUC),
Cornell (\NumCornell),
KTH Sweden (\NumKTH),
Edinburgh UK (\NumEdinburgh).
US universities contribute \NumUSCourses{} courses; international universities contribute \NumIntlCourses{} courses.

\paragraph{What data does each instance consist of?}
Each \emph{course} record contains: \texttt{institution}, \texttt{course\_id}, \texttt{title}, \texttt{description}, \texttt{department}, and \texttt{level} (where available from the source catalog).
Each \emph{occupation} record contains: \texttt{occupation\_id}, \texttt{title}, \texttt{description}, lists of required skills, tasks, and technologies.
Each \emph{mission} record contains: \texttt{institution}, \texttt{statement\_text}, and \texttt{source\_url}.

\paragraph{Is there a label or target associated with each instance?}
No. The benchmark is designed for multi-objective evaluation, not supervised prediction. Evaluation is performed against the Pareto front and benchmark metrics (hypervolume, spacing, coverage).

\paragraph{Is any information missing from individual instances?}
Some fields are unavailable for certain sources: course \texttt{level} is not available from all university catalogs; \texttt{department} naming conventions vary across institutions.
International course descriptions (KTH, Edinburgh) are in English but may reflect translation artifacts from the original Swedish or course-catalog conventions.

\paragraph{Are there any errors, sources of noise, or redundancies?}
Near-duplicate courses (e.g., cross-listed courses) were removed during deduplication. Some course descriptions are short or formulaic (e.g., independent study entries); a minimum text-length filter was applied. O*NET profiles are professionally curated by the US DOL and contain minimal noise.

\paragraph{Does the dataset contain data that might be considered confidential or sensitive?}
No. All data is drawn from public catalogs, public APIs, and public-domain government sources. A basic PII screening step was applied to exclude any inadvertently included names, emails, or phone numbers.

% =====================================================================
\subsection{Collection Process}
% =====================================================================

\paragraph{How was the data associated with each instance acquired?}
Data was acquired from publicly accessible sources:
\begin{itemize}
  \item \textbf{MIT}: MIT OpenCourseWare (OCW) public catalog.
  \item \textbf{UC Berkeley}: UC Berkeley online course catalog.
  \item \textbf{Stanford}: Stanford ExploreCourses public API.
  \item \textbf{UIUC}: UIUC Course Information Suite (CIS) public API.
  \item \textbf{Cornell}: Cornell Class Roster public API.
  \item \textbf{KTH}: KTH KOPPS course syllabus public API (Sweden).
  \item \textbf{Edinburgh}: University of Edinburgh Degree Regulations and Programmes of Study (DRPS) public catalog (UK).
  \item \textbf{O*NET}: US Department of Labor O*NET public database and API.
  \item \textbf{College Scorecard}: US Department of Education College Scorecard public API.
\end{itemize}

\paragraph{What was the collection period?}
Data was collected between 2024 and 2025. O*NET time-travel versions (2022--2025) are included for market-drift analysis.

\paragraph{Were any ethical review processes conducted?}
All sources are public and open-access. No human subjects are involved; no IRB review was required. No personally identifiable information about students, instructors, or workers is collected or included.

\paragraph{Was consent obtained?}
Not applicable. All data consists of publicly published institutional artifacts (course catalogs, government databases).

% =====================================================================
\subsection{Preprocessing, Cleaning, and Labeling}
% =====================================================================

\paragraph{What preprocessing and cleaning was done?}
\begin{enumerate}
  \item \textbf{Schema validation}: all records validated against a required schema (\texttt{institution}, \texttt{course\_id}, \texttt{title} are mandatory fields).
  \item \textbf{Deduplication}: near-identical artifacts within each source were removed.
  \item \textbf{Minimum text length}: artifacts with extremely short or empty descriptions were filtered out.
  \item \textbf{PII screening}: basic screening to exclude names, email addresses, and phone numbers.
  \item \textbf{Provenance tracking}: for every artifact, we record a stable identifier, source URL, snapshot timestamp, content hash of the raw text, and parsing/extraction rules.
\end{enumerate}

\paragraph{Was the raw data saved in addition to the preprocessed data?}
Yes. Provenance manifests link each processed artifact to its raw source. Content hashes allow verification of data integrity.

\paragraph{Is the software used to preprocess the data available?}
Yes. Ingestion scripts and validation code are included in the released repository.

% =====================================================================
\subsection{Uses}
% =====================================================================

\paragraph{What tasks is the dataset intended for?}
The benchmark defines three evaluation tasks:
\begin{itemize}
  \item \textbf{Task A: Multi-Stakeholder Course Ranking.} Given course text and stakeholder target corpora, produce an ordered list of course IDs. Metrics: coverage, diversity, objective balance.
  \item \textbf{Task B: Stakeholder Trade-off Discovery.} Identify a set of courses representing distinct stakeholder trade-offs. Metrics: hypervolume, spacing, spread, Pareto precision, dominance resistance.
  \item \textbf{Task C: Robust Recommendation.} Produce a primary submission plus perturbation-run submissions. Metrics: bootstrap Jaccard, HV retention, rank stability, core stable fraction.
\end{itemize}

\paragraph{What tasks should the dataset \emph{not} be used for?}
The dataset is \emph{not} intended for:
\begin{itemize}
  \item Knowledge tracing or student performance prediction (no learner interaction data is included).
  \item Grade prediction or assessment modeling.
  \item Learner interaction modeling or tutoring system evaluation.
  \item Making enrollment decisions without human oversight.
\end{itemize}

\paragraph{Are there tasks for which the dataset should not be used?}
The dataset should not be used to rank individual students or to automate high-stakes enrollment/admissions decisions. It is a benchmark for evaluating multi-stakeholder trade-off methods, not an operational decision system.

% =====================================================================
\subsection{Distribution}
% =====================================================================

\paragraph{How will the dataset be distributed?}
Upon acceptance, the dataset will be distributed via:
\begin{itemize}
  \item \textbf{Zenodo} with a persistent DOI for long-term archival.
  \item \textbf{HuggingFace Datasets} for programmatic access.
\end{itemize}

\paragraph{When will the dataset be released?}
The dataset will be released upon paper acceptance.

\paragraph{What license will the dataset be distributed under?}
\textbf{CC BY-NC-SA 4.0}. This reflects the most restrictive component license (MIT OCW content is licensed under CC BY-NC-SA 4.0). O*NET data is public domain (US DOL). College Scorecard data is ODC-By (US DOE). University catalog data is from public sources.

\paragraph{Are there any fees or access restrictions?}
No. The dataset, evaluation script (\texttt{benchmark\_eval.py}, requiring only NumPy and the Python standard library), and ingestion code are freely available.

% =====================================================================
\subsection{Maintenance}
% =====================================================================

\paragraph{Who is supporting/hosting/maintaining the dataset?}
The dataset is maintained by the paper authors. Versioned snapshots are released; the current version is \texttt{unified\_v4}.

\paragraph{How can the owner/curator/manager of the dataset be contacted?}
Contact information will be provided upon de-anonymization.

\paragraph{Will the dataset be updated?}
The dataset follows a versioned-snapshot model. The current release (\texttt{unified\_v4}) includes courses from \NumUniversities{} universities. O*NET time-travel versions (2022--2025) are included to support market-drift analysis. Future snapshots may add universities, update course catalogs, or incorporate newer O*NET releases; each snapshot will receive a new version identifier and DOI.

\paragraph{Will older versions continue to be available?}
Yes. Each versioned snapshot is archived on Zenodo with its own DOI. Previous versions will remain accessible for reproducibility.

\paragraph{How will the dataset be retired?}
If a version is found to contain errors, a corrected version will be released and the errata will be documented. No version will be silently removed.
 from main paper appendix.

\section{Datasheet for Datasets}
\label{sec:datasheet}

This datasheet follows the template of Gebru et al.~\cite{gebru2021datasheets}
and documents the ADS cross-national educational benchmark released with this paper.

% =====================================================================
\subsection{Motivation}
% =====================================================================

\paragraph{Why was the dataset created?}
No existing benchmark combines multi-institutional course artifacts with structured labor-market taxonomies and standardized multi-stakeholder evaluation tasks.
Learner-interaction datasets (e.g., KDD Cup 2010, EdNet, OULAD) focus on knowledge tracing at a single institution; MOOC-catalog datasets (MOOCCube, MOOCCubeX) cover multiple institutions but lack labor-market integration.
ADS was created to fill this gap: a cross-national, multi-objective educational benchmark for evaluating trade-offs among learner, institutional, and labor-market stakeholders.

\paragraph{Who created the dataset and on behalf of which entity?}
The dataset was created by the paper authors as an academic research contribution.

\paragraph{Who funded the creation of the dataset?}
No external funding was received for the creation of this dataset.
All data sources are publicly available and were collected using public APIs and catalogs.

% =====================================================================
\subsection{Composition}
% =====================================================================

\paragraph{What do the instances represent?}
The benchmark contains three types of primary text artifacts and two types of structured records:
\begin{itemize}
  \item \textbf{Courses} (\NumCourses): educational course descriptions from \NumUniversities{} universities across \NumCountries{} countries (USA, Sweden, UK).
  \item \textbf{O*NET occupation profiles} (\NumOccupations): occupation descriptions including required skills, tasks, and technologies from the US Department of Labor.
  \item \textbf{Mission statements} (\NumMissionsExpanded): university mission and vision texts.
  \item \textbf{College Scorecard programs} (\NumScorecardPrograms): program-level outcome descriptors with median earnings data (US DOE).
  \item \textbf{Structured attribute/link records} (\NumAttributeRecords): O*NET skills, tasks, technologies, and CIP-SOC crosswalk links.
\end{itemize}

\paragraph{How many instances are there?}
The benchmark contains \NumPrimaryArtifacts{} primary text artifacts and \NumProcessedRows{} total processed rows.
Course counts by university:
MIT OCW (\NumMIT),
UC Berkeley (\NumUCB),
Stanford (\NumStanford),
UIUC (\NumUIUC),
Cornell (\NumCornell),
KTH Sweden (\NumKTH),
Edinburgh UK (\NumEdinburgh).
US universities contribute \NumUSCourses{} courses; international universities contribute \NumIntlCourses{} courses.

\paragraph{What data does each instance consist of?}
Each \emph{course} record contains: \texttt{institution}, \texttt{course\_id}, \texttt{title}, \texttt{description}, \texttt{department}, and \texttt{level} (where available from the source catalog).
Each \emph{occupation} record contains: \texttt{occupation\_id}, \texttt{title}, \texttt{description}, lists of required skills, tasks, and technologies.
Each \emph{mission} record contains: \texttt{institution}, \texttt{statement\_text}, and \texttt{source\_url}.

\paragraph{Is there a label or target associated with each instance?}
No. The benchmark is designed for multi-objective evaluation, not supervised prediction. Evaluation is performed against the Pareto front and benchmark metrics (hypervolume, spacing, coverage).

\paragraph{Is any information missing from individual instances?}
Some fields are unavailable for certain sources: course \texttt{level} is not available from all university catalogs; \texttt{department} naming conventions vary across institutions.
International course descriptions (KTH, Edinburgh) are in English but may reflect translation artifacts from the original Swedish or course-catalog conventions.

\paragraph{Are there any errors, sources of noise, or redundancies?}
Near-duplicate courses (e.g., cross-listed courses) were removed during deduplication. Some course descriptions are short or formulaic (e.g., independent study entries); a minimum text-length filter was applied. O*NET profiles are professionally curated by the US DOL and contain minimal noise.

\paragraph{Does the dataset contain data that might be considered confidential or sensitive?}
No. All data is drawn from public catalogs, public APIs, and public-domain government sources. A basic PII screening step was applied to exclude any inadvertently included names, emails, or phone numbers.

% =====================================================================
\subsection{Collection Process}
% =====================================================================

\paragraph{How was the data associated with each instance acquired?}
Data was acquired from publicly accessible sources:
\begin{itemize}
  \item \textbf{MIT}: MIT OpenCourseWare (OCW) public catalog.
  \item \textbf{UC Berkeley}: UC Berkeley online course catalog.
  \item \textbf{Stanford}: Stanford ExploreCourses public API.
  \item \textbf{UIUC}: UIUC Course Information Suite (CIS) public API.
  \item \textbf{Cornell}: Cornell Class Roster public API.
  \item \textbf{KTH}: KTH KOPPS course syllabus public API (Sweden).
  \item \textbf{Edinburgh}: University of Edinburgh Degree Regulations and Programmes of Study (DRPS) public catalog (UK).
  \item \textbf{O*NET}: US Department of Labor O*NET public database and API.
  \item \textbf{College Scorecard}: US Department of Education College Scorecard public API.
\end{itemize}

\paragraph{What was the collection period?}
Data was collected between 2024 and 2025. O*NET time-travel versions (2022--2025) are included for market-drift analysis.

\paragraph{Were any ethical review processes conducted?}
All sources are public and open-access. No human subjects are involved; no IRB review was required. No personally identifiable information about students, instructors, or workers is collected or included.

\paragraph{Was consent obtained?}
Not applicable. All data consists of publicly published institutional artifacts (course catalogs, government databases).

% =====================================================================
\subsection{Preprocessing, Cleaning, and Labeling}
% =====================================================================

\paragraph{What preprocessing and cleaning was done?}
\begin{enumerate}
  \item \textbf{Schema validation}: all records validated against a required schema (\texttt{institution}, \texttt{course\_id}, \texttt{title} are mandatory fields).
  \item \textbf{Deduplication}: near-identical artifacts within each source were removed.
  \item \textbf{Minimum text length}: artifacts with extremely short or empty descriptions were filtered out.
  \item \textbf{PII screening}: basic screening to exclude names, email addresses, and phone numbers.
  \item \textbf{Provenance tracking}: for every artifact, we record a stable identifier, source URL, snapshot timestamp, content hash of the raw text, and parsing/extraction rules.
\end{enumerate}

\paragraph{Was the raw data saved in addition to the preprocessed data?}
Yes. Provenance manifests link each processed artifact to its raw source. Content hashes allow verification of data integrity.

\paragraph{Is the software used to preprocess the data available?}
Yes. Ingestion scripts and validation code are included in the released repository.

% =====================================================================
\subsection{Uses}
% =====================================================================

\paragraph{What tasks is the dataset intended for?}
The benchmark defines three evaluation tasks:
\begin{itemize}
  \item \textbf{Task A: Multi-Stakeholder Course Ranking.} Given course text and stakeholder target corpora, produce an ordered list of course IDs. Metrics: coverage, diversity, objective balance.
  \item \textbf{Task B: Stakeholder Trade-off Discovery.} Identify a set of courses representing distinct stakeholder trade-offs. Metrics: hypervolume, spacing, spread, Pareto precision, dominance resistance.
  \item \textbf{Task C: Robust Recommendation.} Produce a primary submission plus perturbation-run submissions. Metrics: bootstrap Jaccard, HV retention, rank stability, core stable fraction.
\end{itemize}

\paragraph{What tasks should the dataset \emph{not} be used for?}
The dataset is \emph{not} intended for:
\begin{itemize}
  \item Knowledge tracing or student performance prediction (no learner interaction data is included).
  \item Grade prediction or assessment modeling.
  \item Learner interaction modeling or tutoring system evaluation.
  \item Making enrollment decisions without human oversight.
\end{itemize}

\paragraph{Are there tasks for which the dataset should not be used?}
The dataset should not be used to rank individual students or to automate high-stakes enrollment/admissions decisions. It is a benchmark for evaluating multi-stakeholder trade-off methods, not an operational decision system.

% =====================================================================
\subsection{Distribution}
% =====================================================================

\paragraph{How will the dataset be distributed?}
Upon acceptance, the dataset will be distributed via:
\begin{itemize}
  \item \textbf{Zenodo} with a persistent DOI for long-term archival.
  \item \textbf{HuggingFace Datasets} for programmatic access.
\end{itemize}

\paragraph{When will the dataset be released?}
The dataset will be released upon paper acceptance.

\paragraph{What license will the dataset be distributed under?}
\textbf{CC BY-NC-SA 4.0}. This reflects the most restrictive component license (MIT OCW content is licensed under CC BY-NC-SA 4.0). O*NET data is public domain (US DOL). College Scorecard data is ODC-By (US DOE). University catalog data is from public sources.

\paragraph{Are there any fees or access restrictions?}
No. The dataset, evaluation script (\texttt{benchmark\_eval.py}, requiring only NumPy and the Python standard library), and ingestion code are freely available.

% =====================================================================
\subsection{Maintenance}
% =====================================================================

\paragraph{Who is supporting/hosting/maintaining the dataset?}
The dataset is maintained by the paper authors. Versioned snapshots are released; the current version is \texttt{unified\_v4}.

\paragraph{How can the owner/curator/manager of the dataset be contacted?}
Contact information will be provided upon de-anonymization.

\paragraph{Will the dataset be updated?}
The dataset follows a versioned-snapshot model. The current release (\texttt{unified\_v4}) includes courses from \NumUniversities{} universities. O*NET time-travel versions (2022--2025) are included to support market-drift analysis. Future snapshots may add universities, update course catalogs, or incorporate newer O*NET releases; each snapshot will receive a new version identifier and DOI.

\paragraph{Will older versions continue to be available?}
Yes. Each versioned snapshot is archived on Zenodo with its own DOI. Previous versions will remain accessible for reproducibility.

\paragraph{How will the dataset be retired?}
If a version is found to contain errors, a corrected version will be released and the errata will be documented. No version will be silently removed.
 from main paper appendix.

\section{Datasheet for Datasets}
\label{sec:datasheet}

This datasheet follows the template of Gebru et al.~\cite{gebru2021datasheets}
and documents the ADS cross-national educational benchmark released with this paper.

% =====================================================================
\subsection{Motivation}
% =====================================================================

\paragraph{Why was the dataset created?}
No existing benchmark combines multi-institutional course artifacts with structured labor-market taxonomies and standardized multi-stakeholder evaluation tasks.
Learner-interaction datasets (e.g., KDD Cup 2010, EdNet, OULAD) focus on knowledge tracing at a single institution; MOOC-catalog datasets (MOOCCube, MOOCCubeX) cover multiple institutions but lack labor-market integration.
ADS was created to fill this gap: a cross-national, multi-objective educational benchmark for evaluating trade-offs among learner, institutional, and labor-market stakeholders.

\paragraph{Who created the dataset and on behalf of which entity?}
The dataset was created by the paper authors as an academic research contribution.

\paragraph{Who funded the creation of the dataset?}
No external funding was received for the creation of this dataset.
All data sources are publicly available and were collected using public APIs and catalogs.

% =====================================================================
\subsection{Composition}
% =====================================================================

\paragraph{What do the instances represent?}
The benchmark contains three types of primary text artifacts and two types of structured records:
\begin{itemize}
  \item \textbf{Courses} (\NumCourses): educational course descriptions from \NumUniversities{} universities across \NumCountries{} countries (USA, Sweden, UK).
  \item \textbf{O*NET occupation profiles} (\NumOccupations): occupation descriptions including required skills, tasks, and technologies from the US Department of Labor.
  \item \textbf{Mission statements} (\NumMissionsExpanded): university mission and vision texts.
  \item \textbf{College Scorecard programs} (\NumScorecardPrograms): program-level outcome descriptors with median earnings data (US DOE).
  \item \textbf{Structured attribute/link records} (\NumAttributeRecords): O*NET skills, tasks, technologies, and CIP-SOC crosswalk links.
\end{itemize}

\paragraph{How many instances are there?}
The benchmark contains \NumPrimaryArtifacts{} primary text artifacts and \NumProcessedRows{} total processed rows.
Course counts by university:
MIT OCW (\NumMIT),
UC Berkeley (\NumUCB),
Stanford (\NumStanford),
UIUC (\NumUIUC),
Cornell (\NumCornell),
KTH Sweden (\NumKTH),
Edinburgh UK (\NumEdinburgh).
US universities contribute \NumUSCourses{} courses; international universities contribute \NumIntlCourses{} courses.

\paragraph{What data does each instance consist of?}
Each \emph{course} record contains: \texttt{institution}, \texttt{course\_id}, \texttt{title}, \texttt{description}, \texttt{department}, and \texttt{level} (where available from the source catalog).
Each \emph{occupation} record contains: \texttt{occupation\_id}, \texttt{title}, \texttt{description}, lists of required skills, tasks, and technologies.
Each \emph{mission} record contains: \texttt{institution}, \texttt{statement\_text}, and \texttt{source\_url}.

\paragraph{Is there a label or target associated with each instance?}
No. The benchmark is designed for multi-objective evaluation, not supervised prediction. Evaluation is performed against the Pareto front and benchmark metrics (hypervolume, spacing, coverage).

\paragraph{Is any information missing from individual instances?}
Some fields are unavailable for certain sources: course \texttt{level} is not available from all university catalogs; \texttt{department} naming conventions vary across institutions.
International course descriptions (KTH, Edinburgh) are in English but may reflect translation artifacts from the original Swedish or course-catalog conventions.

\paragraph{Are there any errors, sources of noise, or redundancies?}
Near-duplicate courses (e.g., cross-listed courses) were removed during deduplication. Some course descriptions are short or formulaic (e.g., independent study entries); a minimum text-length filter was applied. O*NET profiles are professionally curated by the US DOL and contain minimal noise.

\paragraph{Does the dataset contain data that might be considered confidential or sensitive?}
No. All data is drawn from public catalogs, public APIs, and public-domain government sources. A basic PII screening step was applied to exclude any inadvertently included names, emails, or phone numbers.

% =====================================================================
\subsection{Collection Process}
% =====================================================================

\paragraph{How was the data associated with each instance acquired?}
Data was acquired from publicly accessible sources:
\begin{itemize}
  \item \textbf{MIT}: MIT OpenCourseWare (OCW) public catalog.
  \item \textbf{UC Berkeley}: UC Berkeley online course catalog.
  \item \textbf{Stanford}: Stanford ExploreCourses public API.
  \item \textbf{UIUC}: UIUC Course Information Suite (CIS) public API.
  \item \textbf{Cornell}: Cornell Class Roster public API.
  \item \textbf{KTH}: KTH KOPPS course syllabus public API (Sweden).
  \item \textbf{Edinburgh}: University of Edinburgh Degree Regulations and Programmes of Study (DRPS) public catalog (UK).
  \item \textbf{O*NET}: US Department of Labor O*NET public database and API.
  \item \textbf{College Scorecard}: US Department of Education College Scorecard public API.
\end{itemize}

\paragraph{What was the collection period?}
Data was collected between 2024 and 2025. O*NET time-travel versions (2022--2025) are included for market-drift analysis.

\paragraph{Were any ethical review processes conducted?}
All sources are public and open-access. No human subjects are involved; no IRB review was required. No personally identifiable information about students, instructors, or workers is collected or included.

\paragraph{Was consent obtained?}
Not applicable. All data consists of publicly published institutional artifacts (course catalogs, government databases).

% =====================================================================
\subsection{Preprocessing, Cleaning, and Labeling}
% =====================================================================

\paragraph{What preprocessing and cleaning was done?}
\begin{enumerate}
  \item \textbf{Schema validation}: all records validated against a required schema (\texttt{institution}, \texttt{course\_id}, \texttt{title} are mandatory fields).
  \item \textbf{Deduplication}: near-identical artifacts within each source were removed.
  \item \textbf{Minimum text length}: artifacts with extremely short or empty descriptions were filtered out.
  \item \textbf{PII screening}: basic screening to exclude names, email addresses, and phone numbers.
  \item \textbf{Provenance tracking}: for every artifact, we record a stable identifier, source URL, snapshot timestamp, content hash of the raw text, and parsing/extraction rules.
\end{enumerate}

\paragraph{Was the raw data saved in addition to the preprocessed data?}
Yes. Provenance manifests link each processed artifact to its raw source. Content hashes allow verification of data integrity.

\paragraph{Is the software used to preprocess the data available?}
Yes. Ingestion scripts and validation code are included in the released repository.

% =====================================================================
\subsection{Uses}
% =====================================================================

\paragraph{What tasks is the dataset intended for?}
The benchmark defines three evaluation tasks:
\begin{itemize}
  \item \textbf{Task A: Multi-Stakeholder Course Ranking.} Given course text and stakeholder target corpora, produce an ordered list of course IDs. Metrics: coverage, diversity, objective balance.
  \item \textbf{Task B: Stakeholder Trade-off Discovery.} Identify a set of courses representing distinct stakeholder trade-offs. Metrics: hypervolume, spacing, spread, Pareto precision, dominance resistance.
  \item \textbf{Task C: Robust Recommendation.} Produce a primary submission plus perturbation-run submissions. Metrics: bootstrap Jaccard, HV retention, rank stability, core stable fraction.
\end{itemize}

\paragraph{What tasks should the dataset \emph{not} be used for?}
The dataset is \emph{not} intended for:
\begin{itemize}
  \item Knowledge tracing or student performance prediction (no learner interaction data is included).
  \item Grade prediction or assessment modeling.
  \item Learner interaction modeling or tutoring system evaluation.
  \item Making enrollment decisions without human oversight.
\end{itemize}

\paragraph{Are there tasks for which the dataset should not be used?}
The dataset should not be used to rank individual students or to automate high-stakes enrollment/admissions decisions. It is a benchmark for evaluating multi-stakeholder trade-off methods, not an operational decision system.

% =====================================================================
\subsection{Distribution}
% =====================================================================

\paragraph{How will the dataset be distributed?}
Upon acceptance, the dataset will be distributed via:
\begin{itemize}
  \item \textbf{Zenodo} with a persistent DOI for long-term archival.
  \item \textbf{HuggingFace Datasets} for programmatic access.
\end{itemize}

\paragraph{When will the dataset be released?}
The dataset will be released upon paper acceptance.

\paragraph{What license will the dataset be distributed under?}
\textbf{CC BY-NC-SA 4.0}. This reflects the most restrictive component license (MIT OCW content is licensed under CC BY-NC-SA 4.0). O*NET data is public domain (US DOL). College Scorecard data is ODC-By (US DOE). University catalog data is from public sources.

\paragraph{Are there any fees or access restrictions?}
No. The dataset, evaluation script (\texttt{benchmark\_eval.py}, requiring only NumPy and the Python standard library), and ingestion code are freely available.

% =====================================================================
\subsection{Maintenance}
% =====================================================================

\paragraph{Who is supporting/hosting/maintaining the dataset?}
The dataset is maintained by the paper authors. Versioned snapshots are released; the current version is \texttt{unified\_v4}.

\paragraph{How can the owner/curator/manager of the dataset be contacted?}
Contact information will be provided upon de-anonymization.

\paragraph{Will the dataset be updated?}
The dataset follows a versioned-snapshot model. The current release (\texttt{unified\_v4}) includes courses from \NumUniversities{} universities. O*NET time-travel versions (2022--2025) are included to support market-drift analysis. Future snapshots may add universities, update course catalogs, or incorporate newer O*NET releases; each snapshot will receive a new version identifier and DOI.

\paragraph{Will older versions continue to be available?}
Yes. Each versioned snapshot is archived on Zenodo with its own DOI. Previous versions will remain accessible for reproducibility.

\paragraph{How will the dataset be retired?}
If a version is found to contain errors, a corrected version will be released and the errata will be documented. No version will be silently removed.
